\Askhsh{14805}{3}
\begin{ylh}{2}
    \item 2.1. Οι Πράξεις και οι Ιδιότητές τους
    \item 2.3. Απόλυτη Τιμή Πραγματικού Αριθμού
    \item 2.4. Ρίζες Πραγματικών Αριθμών
\end{ylh}
Έστω α ένας πραγματικός αριθμός, για τον οποίο ισχύει $a = |3\sqrt{2}-4| + |2\sqrt{2}-2|$.

\begin{alist}
\item Να αποδείξετε ότι $a = 2$. (Θεωρήστε ότι $\sqrt{2} \approx 1,41$) \monades{10}
\item Με τη βοήθεια του ερωτήματος (α) να αποδείξετε ότι $a^3 = 2a$. \monades{5}
\item Να βρείτε την αριθμητική τιμή της παράστασης $A = a^3 + (a-1)^2$. \monades{10}
\end{alist}

